% Source of appendix.tex.  build_appendix.py fills each %%...%% placeholder from
% ../greedy_intersection.tex.

\section{Proofs for Section~\texorpdfstring{\ref{sec:model}}{2}}
\label{app:model}

We use the following classical facts~\cite{oxley11,schrijver03}.

%%ENV:lem:minbases%%

Explicitly, $M^{u}$ is the direct sum, over the values $\ell$ of $u$, of the
matroids $(M|\{u\le\ell\})/\{u<\ell\}$ on the level sets $\{u=\ell\}$, so its
rank function is read from that of $M$: if $\ell_1<\cdots<\ell_q$ are the values
of $u$ and $U_i=\{u\le\ell_i\}$, $U_0=\varnothing$, then
\begin{equation}
  r_{M^{u}}(A)=\sum_{i=1}^{q}\Bigl(r\bigl(U_{i-1}\cup(A\cap\{u=\ell_i\})\bigr)-r(U_{i-1})\Bigr).
  \label{eq:rank-mu}
\end{equation}

The two expressions in Definition~\ref{def:ti} agree, as $G_t(D)=G_{s_i}(D)$ for
$t\in[s_i,s_{i+1})$.  If $F(E)=0$, the upper limit $s_m=\max_eD_e$ can be replaced
by $\infty$, as $G_t(D)=E$ for $t\ge s_m$.

%%PROOFREF:thm:compiler%%

%%ENV:cor:integer%%

%%PROOF:cor:integer%%

%%PROOFREF:cor:levels%%

\section{Proofs for Section~\texorpdfstring{\ref{sec:greedy}}{3}}
\label{app:greedy}

%%PROOFREF:lem:greedy%%

%%PROOFOF:thm:A%%

%%PROOFOF:thm:B%%

%%ENV:cor:equiv%%

%%PROOFREF:cor:classes%%

A Turing machine that runs in time $T$ on inputs of a fixed length is simulated
by a Boolean circuit of size $O(T\log T)$~\cite{pippenger-fischer79}.  So every
matroid whose rank function is computable in time polynomial in $m$ and in the length of a fixed
description has a Boolean rank circuit, and by Theorem~\ref{thm:A} a
$(\max,+,-)$-circuit for $h_M$, of size polynomial in these two quantities.

%%ENV:cor:count%%

%%PROOF:cor:count%%

\section{Proofs for Section~\texorpdfstring{\ref{sec:rainbow}}{4}}
\label{app:rainbow}

\subsection{Setting and the Lagrangian}
\label{app:lagr}

The definitions of Section~\ref{sec:lagr} and Lemmas~\ref{lem:avoidable}
and~\ref{lem:augment} apply verbatim to every partition of a finite ground set
into any number of blocks; Appendix~\ref{app:two} uses them in this form.
Rainbow independent sets of $M^{v}$ extend to minimum bases of $v$, and a minimum
base meeting $j$ blocks contains a rainbow set of size $j$; so $\nu_0$ is the
largest number of blocks met by a minimum base of $v$, and $S_0$ is the set of
blocks missed by some minimum base meeting $\nu_0$ blocks.

%%PROOFREF:lem:duality%%

%%ENV:prop:star%%

%%PROOF:prop:star%%

\subsection{Possibly Optimal Bases and Two Obstacles}
\label{app:pos}

%%PROOFREF:prop:pos%%

%%ENV:lem:path%%

%%PROOF:lem:path%%

\begin{proof}[Proof of Proposition~\ref{prop:obstacles}(1)]
%%PROOFBODY:prop:notmatroid%%
\end{proof}

\begin{proof}[Proof of Proposition~\ref{prop:obstacles}(2)]
%%PROOFBODY:prop:nphard%%
\end{proof}

\subsection{The Update}
\label{app:update}

\begin{proof}[Proof of Lemma~\ref{lem:updint}]
By (F2) independent sets of $N(t)$ stay independent in $N(t')$ for $t\le t'$, and
$\nu(N(0))=\nu_0\ge r$ by (F1), so
\[
  r\le\nu(N(0))\le\nu(N(t))\le\nu(N(t'))\qquad(t\le t').
\]
Moreover $\nu(N(t))=k>r$ for large $t$ by (F4), and by (F5) both sets of levels in
Definition~\ref{def:update} are unions of intervals $[c_i,c_{i+1})$ and
$[c_q,\infty)$, so the minima exist.  For $t<\tau$ we have $\nu(N(t))=r$; a base of
$N(t)$ meeting exactly $r$ blocks contains a rainbow independent set of size $r$
missing the same blocks, and a rainbow independent set of size $r$ extends to a
base of $N(t)$ meeting no further block.  With (F2) this gives
$b\in\Av(N(t))\iff\lambda_b\le t$ for $t<\tau$, and as
$\mathbf 1[\nu(N(t))\le r]=1$ exactly for $t<\tau$,
\[
  \int_0^{\infty}\mathbf 1\bigl[\nu(N(t))\le r\bigr]\;\mathbf 1\bigl[b\in\Av(N(t))\bigr]\,dt
  =\int_0^{\tau}\mathbf 1[\lambda_b\le t]\,dt=(\tau-\lambda_b)_+=a_b.\qedhere
\]
\end{proof}

%%ENV:lem:avoidable%%

%%PROOF:lem:avoidable%%

%%PROOFREF:thm:validity%%

\subsection{Single Linkage on Possibly Optimal Elements}
\label{app:posl}

Both inequalities in Definition~\ref{def:posl} are needed: with $v_g\le v_f-t$ in
$\PO_t$, (F1) fails for $U_{1,3}$ and $v=(0,0,1)$, and with $|v_x-v_y|<t$ in the
linkage, the two elements of value $0$ there form two clusters at scale $0$ that
are not ordered.

%%PROOFREF:prop:posl%%

\subsection{Breakpoints}
\label{app:events}

%%PROOFREF:lem:events%%

\subsection{Termination}
\label{app:term}

The ties in $\Pi(v)$ are needed: on the cycle matroid of $K_4$ with the edges
$01$, $02$, $03$, $12$, $13$, $23$ in the blocks $R_1$, $R_1$, $R_1$, $R_1$, $R_2$,
$R_3$ and $v=(0,0,3,1,2,2)$, where $\nu_0=2$, $S_0=\{2,3\}$, $\tau=1$ and
$a=(0,1,1)$, the update adds the pairs $(13,12)$ and $(23,12)$ to $\Pi$, both with
equality; with strict inequality, item 3 of Theorem~\ref{thm:potential} fails.

%%PROOFREF:thm:potential%%

%%CROSSING%%

%%STAGE:lemma%%

%%STAGE:theorem%%

%%PROOFREF:cor:rankone%%

\subsection{The Circuit}
\label{app:circuit}

%%PROOFREF:lem:augment%%

%%PROOFREF:lem:observer%%

%%PROOFREF:lem:oneupdate%%

%%PROOFOF:thm:C%%

\section{Proofs for Section~\texorpdfstring{\ref{sec:two}}{5}}
\label{app:two}

%%PROOFREF:lem:doubling%%

%%PROOFOF:thm:D%%

The circuit of Lemma~\ref{lem:augment} computes the largest size of a rainbow
independent set, that is, it solves unweighted matroid intersection with a
partition matroid.  Transversal matroids have rank functions of this kind, and
gammoids are restrictions of their duals, so Theorem~\ref{thm:A} applies to
both.

%%ENV:cor:transversal%%

%%PROOF:cor:transversal%%

\section{Formal Verification and Computational Checks}
\label{app:checks}

\paragraph{Formal proofs.}  A Lean~4 development with Mathlib, both at version
4.32.0, proves the results listed below with the constants stated in the paper, for the
matroids of Mathlib's \texttt{Matroid} on $\{0,\ldots,m-1\}$.

For Sections~\ref{sec:model} and~\ref{sec:greedy} it proves
Lemma~\ref{lem:minbases}, the equality of the two expressions in
Definition~\ref{def:ti}, Theorem~\ref{thm:compiler},
Corollaries~\ref{cor:integer} and~\ref{cor:levels}, Lemma~\ref{lem:greedy},
Theorems~\ref{thm:A} and~\ref{thm:B}, and Corollary~\ref{cor:classes}.  For
Corollary~\ref{cor:count} it proves the count of Step~1 and the conclusion in an
explicit form: for $m\ge128$, at most a fraction $2^{-\lfloor2^{m}/m^{3}\rfloor}$
of the matroids on $m$ labeled elements have a $(\max,+,-)$-circuit of size less
than $2^{m/2}/m^{2}$ that computes $h_M$ on $\{0,1\}^{E}$.

For Sections~\ref{sec:rainbow} and~\ref{sec:two} it proves
Lemmas~\ref{lem:duality}, \ref{lem:path}, \ref{lem:updint}, \ref{lem:avoidable},
\ref{lem:events}, \ref{lem:crossing}, \ref{lem:augment}, \ref{lem:observer},
\ref{lem:oneupdate} and~\ref{lem:doubling}, Propositions~\ref{prop:star},
\ref{prop:pos} and~\ref{prop:posl}, Proposition~\ref{prop:obstacles}(1),
Theorems~\ref{thm:validity}, \ref{thm:potential} and~\ref{thm:stagelength},
Corollary~\ref{cor:rankone}, Theorem~\ref{thm:C} with the size
$3328\,k^{3}m^{4}(s+m)+2$, and $1390\,k^{3}m^{3}(s+m)+2$ when every component of
$M$ has rank at most one, Theorem~\ref{thm:D} with the sizes
$3328\cdot1200\,m^{7}(s_1+s_2+m)+2$ for common bases and
$3328\cdot653313\,m^{7}(s_1+s_2+m)+3$ for common independent sets, and
Corollary~\ref{cor:transversal}.  For Proposition~\ref{prop:obstacles}(2) it
proves the reduction, Steps 2 to 4, for every multigraph $H$ and every
$h\ge|V(H)|,|C|$; the complexity classes are not formalized, and the
NP-completeness of the rainbow path problem is cited.
Proposition~\ref{prop:ladder} is not formalized; it is checked by computation
below.  Uniformity, a property of the algorithm that produces the circuits, is
not formalized either.  The following have no Lean statements:
Corollary~\ref{cor:equiv}, which combines Theorems~\ref{thm:A} and~\ref{thm:B};
the remark on simulating computations by Boolean circuits in
Appendix~\ref{app:greedy}; and the two examples opening Appendices~\ref{app:posl}
and~\ref{app:term}.  The main
theorems depend on no axioms beyond the three standard axioms of Lean, and the
build reports no warnings.

%%CHECKS%%

\section{Stage Length of PO-SL}
\label{app:scope}

On the following inputs, the stage length of PO-SL is not bounded in terms of
$r$ and $k$ alone, while the exact family $\Pos_t$, which is not admissible,
needs one update.

%%ENV:prop:ladder%%

%%PROOF:prop:ladder%%
